\scanpage{29}\mspage{V}
\deleted{$]z,x]\subset\mathbb{R}^{2}-\mathbb{C}$, c.q.f.d.}

% slanted
\margin{\uline{visa}}
\uline{Lemme} V. Soit $\mathbb{C}$ un \uline{ensemble convexe} nonvide de $\mathbb{R}^{2}$ et soit $y\in\mathbb{C}$ tel qu'il existe $s,s'\in\mathbb{C}$ avec $y\in ]s,s'[$. Supposons que $y$ est $(\mathbb{R}^{2}-C)$-visible d'un point $x\notin\mathbb{C}$ qui n'est pas colineaire avec $s$ et $s'$. Alors tout point colineaire avec $s$ et $s'$ est \inserted{$(\mathbb{R}^{2}-\mathbb{C})$-}visible de $x$.

\margin{\uline{NB} $s,s'$ \unsure{appartiennent} \unsure{:|à} une $=$ \unsure{arête|arète} de $C$, et si $D$ est la droite qu'elle détermine $x$ et $C$ sont de part et d'autre de cette droite\ldots}

\uline{Demonstration} : Soit $d$ la droite determinée par $s$ et par $s'$, et appelons $\mathbb{S}_{1}$ et $\mathbb{S}_{2}$ les deux demi-espaces ouverts dans lesquels $\mathbb{R}^{2}$ est partagé par $d$.

Comme $x$ n'est pas colineaire avec $s$ et $s'$, \deleted{\unsure{ou|on}} $x\in S_{1}$ ou bien $x\in S_{2}$.

Supposons que $x\in\mathbb{S}_{1}$ et $z\in d$ et prouvons que $]x,z[\subset\mathbb{R}^{2}-\mathbb{C}$. Si $]x,z[\not\subset\mathbb{R}^{2}\deleted{,}-\mathbb{C}$ il doit exister $w\in ]x,z[\cap\mathbb{C}{}^{\inserted{\subset\mathbb{S}_{1}}}$. Comme $w\in\mathbb{S}_{1}$ aussi, les trois points $s_{\deleted{\unsure{1}}},s_{\deleted{\unsure{2}}}',w$ sont affinement independants.

Considerons donc le triangle $T$ de sommets $s,s',w$ (i.e. l'enveloppe convexe de $\{s',s,w\}$). Comme $[x,\deleted{\unsure{z}}y]\cap ]s,s'[\ne\phi$ on a aussi $[x,y\deleted{\unsure{]}}[\cap T\ne\phi$, mais $T\subset\mathbb{C}$ puisque $\{s,s',w\}\subset\mathbb{C}$ ce qui est un absurd \inserted{puisque on suppose $[x,y[\cap\mathbb{C}=\phi$}, cela prouve donc le lemme.

% A short bent pen stroke in the left margin faces the heading.
\margin{\tikz[baseline]{\draw[line cap=round] (0,0) .. controls (-.02,.16) and (0,.21) .. (.08,.23) -- (.34,.36);}}
\uline{Definition} 3. Un polygone $\Gamma$ est \inserted{dit} convexe \inserted{(par abus de langage !)} si l'ensemble sous-jacent à $\Gamma$ est le bord d'un ensemb\uline{le} \uline{co}nvexe (i.e. $\Gamma=\mathbb{C}-\mathring{\mathbb{C}}$, $\mathbb{C}$ convexe)

\margin{$\bullet$}
On peut maintenant én\deleted{n}oncer le :

\uline{Lemme} VI. Si $\Gamma$ est un polygone convexe alors $\Gamma$ est un J-polygone et tout sommet de $\Gamma$ \deleted{\illegible{1}} qui n'est pas colinéaire avec \deleted{\unsure{s}}les deux sommets adjacents ($=$ consecutifs) est un bon sommet (cf. def.2)
% archive sheet number: 28
